%!TEX program = latex
% LaTeX resume using res.cls
\documentclass[margin]{res}
\usepackage{graphicx}
\usepackage{amssymb}
\usepackage{amsmath}
\usepackage{subfigure}
\usepackage{makeidx}
%\usepackage{graphicx}
%\usepackage{helvetica} % uses helvetica postscript font (download helvetica.sty)
%\usepackage{newcent}   % uses new century schoolbook postscript font
\setlength{\textwidth}{5.1in} % set width of text portion
\setlength{\marginparsep}{-80pt}
\begin{document}

% Center the name over the entire width of resume:
 \moveleft.5\hoffset\centerline{\large\bf Lingyan Ran}
% Draw a horizontal line the whole width of resume:
 \moveleft\hoffset\vbox{\hrule width\resumewidth height 1pt}\smallskip

% address begins here
% Again, the address lines must be centered over entire width of resume:
 \moveleft.5\hoffset\centerline{Ph.D Candidate}
 \moveleft.5\hoffset\centerline{School of Computer Science and Technology, 886 mail box}
 \moveleft.5\hoffset\centerline{Chang'an Campus, Northwestern Polytechnical University }
 \moveleft.5\hoffset\centerline{Xi'an, Shaanxi, 710129, China}
 \moveleft.5\hoffset\centerline{Tel: +86-17791834424}
 \moveleft.5\hoffset\centerline{E-mail: lingyanran@gmail.com}
 \moveleft.5\hoffset\centerline{http://saiip-vision.org/LingyanRan/index.htm}
\begin{resume}
\marginpar{\includegraphics[width=0.9in]{photo.eps}}
%\section{OBJECTIVE}  Supported by the China Scholarship Council (CSC) , I am now applying for a position
%                in the field of computer vision as a joint-cultivation PhD candidate, with special
%                interests in binary feature description,
%                image registration, stitching and video stabilization.
\section{SUPERVISOR\\LABORATORY} Supervised by Prof. Yanning Zhang, Dean of the school of Computer Science and Technology\\ From Digital Video Processing Group, Shaanxi Provincial Key Laboratory of Speech \& Image Information Processing

\section{EDUCATION}
                {\sl Concentration: Computer Science and Technology}\\
                {\sl\bf Visiting student}, Stevens Institute of Technology(SIT), Hoboken, USA \hfill Sep. 2013 - Sep. 2015\\
                Northwestern Polytechnical University(NPU), Xi'an, China\\
                {\sl\bf Ph.D candidate,} school of Computer Science and Technology\hfill Sep. 2011 - Present\\
                {\sl\bf Bachelor,} school of Computer Science and Technology \hfill Sep. 2010 - July 2011\\
                {\sl\bf Bachelor,} Honors College \hfill Sep. 2007 - July 2010
                      % \sl will be bold italic in New Century Schoolbook (or
	              % any postscript font) and just slanted in
		      %	Computer Modern (default) font

                %Minor: Management
\section{HONORS  \& \\AWARDS}
                First Level Prize of Excellent Graduate Student of NPU \hfill 2011\\
                First Level Prize of Excellent Graduate Paper for Bachelor Degree\hfill 2010\\
                First Level Prize of Excellent Student of NPU\hfill 2007-2010\\
                Outstanding Position Paper Award from the Model UN Conference, NPU \hfill 2008

\section{ENGLISH PROFICIENCY}
                College English Test (CET) - Band 6

\section{COMPUTER \\ SKILLS} {\sl Programming Language:} C/C++, Matlab, Lua, etc\\
                {\sl Programming Tools:} Microsoft Visual Studio, Matlab\\
                {\sl Programming Library:} OpenCV, Torch, MFC

\section{PROFESSIONAL\\ACTIVITIES}
                Teaching Assistant for the Introduction to Audio, Speech \& Langauge Processing Class\hfill 2012\\
                Volunteer for the Sino-foreign-interchange Workshop on Intelligence Science and Intelligent Data Engineering (ISciDE)\hfill 2011\\
                Volunteer for the Asia-Pacific Signal and Information Processing Association (APSIPA)\hfill 2011\\
                Monitor of the NPU Honorable Model Class\hfill 2007-2010

\section{RESEARCH \\INTERESTS}
                Computer Vision and Pattern Recognition\\
                Deep Learning\\
                Binary Feature Description\\
                Multi-sensor Information Cooperative Processing\\
                Camera Array Synthetic Aperture Imaging\\
                Image Registration and Stitching
                %Video Stabilization, Visual Object Detection, Tracking and Recognition

\section{PUBLICATIONS}
                \textbf{Lingyan Ran}, Yanning Zhang, Gang Hua, CANNET: Context Aware Nonlocal Convolutional Networks for Semantic Image Segmentation, International Conference on Image Processing (ICIP 2015), Quebec, Canada.
                
                TaoYang, Jing Li, Jingyi Yu, Yanning Zhang, Wenguang Ma, Xiaomin Tong, Rui Yu, \textbf{Lingyan Ran}. Multiple-Layer Visibility Propagation-Based Synthetic Aperture Imaging through Occlusion. Sensors, 2015, 15, 18965-18984.
    
                Tao Yang, Yanning Zhang, Jingyi Yu, Jing Li, Wenguang Ma, Xiaomin Tong, Rui Yu, \textbf{Lingyan Ran}. All-In-Focus Synthetic Aperture Imaging, ECCV 2014.

                Tao Yang, Yanning Zhang, Rui Yu, Xiaoqiang Zhang, Ting Chen, \textbf{Lingyan Ran}, Zhengxi Song. Simultaneous
                camera array focus plane estimation and occluded moving object imaging. Image and Vision Computing, 2014

                Tao Yang, Xiaoqiang Zhang, \textbf{Lingyan Ran}, Rui Yu, Runping Xi, {\sl Camera Array SLynthetic Aperture Focusing and
                Fusion based Hidden Object Imaging,} Sino-foreign-interchange Workshop on Intelligence Science and Intelligent
                Data Engineering (ISciDE 2011), Xi’an, China.

\section{RESEARCH \\EXPERIENCE}

                {\bf Invariant Feature Based Automatic Video Stitching in Wide Aerial Surveillance} \hfill (Feb. 2012 - Sep. 2012), NPU\\
                {\sl Project Description: } Algorithms for aligning images and stitching them into seamless photo-mosaics are among the
                oldest and most widely used in computer vision. Image stitching algorithms create the highresolution
                photo-mosaics used to produce today’s digital maps and satellite photos. They also
                come bundled with most digital cameras currently being sold, and can be used to create beautiful
                ultra wide-angle panoramas. Powerful methods such as autostich has been widely used. When it comes to handle the real-time video stitching problem that captured by a camera on a moving platform, those methods seems to perform too much computing. We expected to find an efficient solution for this particular task, using low cost approaches such as Minimum Spanning Tree (MST), etc to perform the relationships between video frames.\\
                {\sl Experience: }
                \begin{itemize} \itemsep -2pt
                \item I presented a fast approach for real-time video sititching using the MST method, its high efficiency was well accepted when dealing with the real-time moving platform video stitching problem.
                \end{itemize}

                %{\sl National Natural Science Foundation of China (NSFC)}
                {\bf Camera Array Synthetic Aperture Focusing and Fusion} \hfill (Sep. 2011 - Jan. 2012), NPU\\
                {\sl Project Description: } Hidden object imaging is challenging problem in the fields
                of computer vision and image processing, and it’s a key step in many
                application fields, include intelligent video surveillance, visual tracking
                and scene understanding. Recently, the camera array synthetic aperture
                imaging has been proved to be a powerful technology for hidden object
                detection, and the state-of-art synthetic imaging method can focus on
                multiple parallel planes so as to achieve seeing hidden through severe
                occlusion. However, due to the depth variation of hidden object’s surface,
                it’s difficult for existing methods to get a complete clear image. Our goal is to find a novel framework for high performance hidden object
                imaging in which we integrate the camera array synthetic aperture
                imaging technique with multi-resolution image fusion method together, and also
                developed a hidden object imaging system with AXIS network camera array. \\
                {\sl Experience: }
                \begin{itemize} \itemsep -2pt %reduce space between items
                %\item In this research I got more detailed information on Synthetic Aperture Imageing theory and got impressed by its advantage in dealing with occlusion problems.
                \item Together we proposed a novel framework adopting several multiple scale image fusion methods to create high resolution and highly detailed images of occluded objects.
                \end{itemize}

                {\bf Real-time Vision-based Autonomous Unmanned Aerial Vehicles(UAV) Landing System}\hfill(June, 2011 - Oct. 2011),NPU\\
                {\sl Project Description: } Vision-based autonomous navigation is one of the hot and difficult spots of current researches in computer vision. A vision system on board an UAV typically arguments a sensor suite that might include Global Positioning System, Inertial Navigation Sensors, laser range finders, a digital compass and sonar. The design of any real-time vision system is a daunting task: It involves a systematic integration of hardware, low level image processing, multi-view geometry and synthesis of real-time controllers. This project aims to solve the problem of vision-based autonomous landing of an UAV, particularly on the technology of ground-based autonomous landing system, without any other sensors or assistances.
                {\sl Experience: }
                \begin{itemize} \itemsep -2pt
                \item I accomplished part of the pre-processing system, including highly accurate calibration of camera intrinsic and extrinsic parameters in wide-area scene, detection 、localization and tracking of an onboard cooperative marker, etc.
                \end{itemize}

                {\bf Moving Object Trajectory Based Infrared and Color Videos Registration} \hfill (Sep. 2010 - June , 2011), NPU\\
                {\sl Project Description: } Video surveillance has been widely used in many fields in recent years and moving target detection and tracking is an important part of it. Also tracking is the basis for further research on behavior recognition. The traditional target detection systems are limited by the sensors they use, so they cannot fulfill the demand in some long-term surveillance scenes. The visible and thermal based multi-sensor tracking system has been paid attention lately. This project focuses on the moving target detection and multi-target tracking algorithm and some other methods like shadow suppression and multi-source image fusion using visible and infrared spectrum videos.\\
                {\sl Experience: }
                \begin{itemize} \itemsep -2pt
                \item I presented a method and a novel criterion to
                register infrared and color (visible) videos. It is a feature point-based
                method that uses top pixel coordinates found after foreground
                detection and tracking to build trajectories which are used to find the transformation
                matrix in both visible and infrared videos.
                \end{itemize}

\section{SELF-\\EVALUATION}
            Great team spirit\\
            Independent technical researching ability\\
            Full of enthusiasm in work\\
            Great communication skills with colleagues and clients

\section{REFEREES}
            [1] Yanning Zhang, Professor, Ph.D.\\
            \hangafter 1
　　        \hangindent 1.5em
　　        \noindent
            Director of Shaanxi Provincial Key Lab of Speech and Information Processing \\
            Dean of School of Computer Science, NPU, Xi’an, China\\
            Phone: 86-29-88431536 (Office); 86-13060399678 (cell)\\
            Email: ynzhang@nwpu.edu.cn \\
            Website: http://www.saiip-vision.org/ynzhang.htm

            [2] Gang Hua, Senior Research Manager, Ph.D.\\
            \hangafter 1
　　        \hangindent 1.5em
　　        \noindent
            Senior Research Manager, MSRA, Beijing, China\\
            Visiting associate professor, Hoboken, NJ, USA\\
            Phone: 86-18410216795(cell)\\
            Email: ganghua@microsoft.com\\
            Website: http://www.cs.stevens.edu/~ghua/

            [3] Tao Yang, Professor, Ph.D.\\
            \hangafter 1
　　        \hangindent 1.5em
　　        \noindent
            Profession of School of Computer Science, NPU, Xi’an, China\\
            Phone: 86-15002919079(cell)\\
            Email: yangtaonwpu@163.com\\
            Website: http://www.saiip-vision.org/tyang/index.html

\hfill Thanks for reading!

\end{resume}
\end{document}




